% ============================================================
%  Sports Manager — Design and Architecture Document
%  CE 216, Izmir University of Economics
%  Compile with: pdflatex report.tex  (or XeLaTeX for full UTF-8)
% ============================================================

\documentclass[12pt, a4paper]{article}

% ── Packages ────────────────────────────────────────────────
\usepackage[utf8]{inputenc}
\usepackage[T1]{fontenc}
\usepackage[english]{babel}
\usepackage{geometry}
\usepackage{graphicx}
\usepackage{booktabs}
\usepackage{longtable}
\usepackage{array}
\usepackage{tabularx}
\usepackage{enumitem}
\usepackage{hyperref}
\usepackage{listings}
\usepackage{xcolor}
\usepackage{titlesec}
\usepackage{fancyhdr}
\usepackage{caption}
\usepackage{float}
\usepackage{parskip}
\usepackage{microtype}
\usepackage{setspace}
\usepackage{wrapfig}
\usepackage{placeins}

% ── Custom column types ───────────────────────────────────────
\newcolumntype{L}[1]{>{\raggedright\arraybackslash}p{#1}}

% ── Page geometry ────────────────────────────────────────────
\geometry{
    top    = 2.5cm,
    bottom = 2.5cm,
    left   = 3.0cm,
    right  = 2.5cm
}

% ── Hyperref setup ───────────────────────────────────────────
\hypersetup{
    colorlinks = true,
    linkcolor  = black,
    urlcolor   = blue,
    citecolor  = black,
    pdftitle   = {Sports Manager – Design and Architecture Document},
}

% ── Section title formatting ─────────────────────────────────
\titleformat{\section}{\large\bfseries}{\thesection.}{0.5em}{}
\titleformat{\subsection}{\normalsize\bfseries}{\thesubsection}{0.5em}{}
\titleformat{\subsubsection}{\normalsize\itshape}{\thesubsubsection}{0.5em}{}

% ── Code listing style ───────────────────────────────────────
\lstset{
    language     = Java,
    basicstyle   = \ttfamily\small,
    keywordstyle = \bfseries\color{blue!70!black},
    commentstyle = \itshape\color{gray},
    stringstyle  = \color{orange!80!black},
    breaklines   = true,
    frame        = single,
    framerule    = 0.4pt,
    rulecolor    = \color{gray!40},
    backgroundcolor = \color{gray!5},
    xleftmargin  = 1em,
    xrightmargin = 1em,
    numbers      = none,
}

% ── Header / footer ──────────────────────────────────────────
\pagestyle{fancy}
\fancyhf{}
\rhead{\small CE 216 – Sports Manager}
\lhead{\small Design and Architecture Document}
\cfoot{\thepage}
\renewcommand{\headrulewidth}{0.4pt}

% ── Table column types ───────────────────────────────────────
\newcolumntype{L}[1]{>{\raggedright\arraybackslash}p{#1}}
\newcolumntype{C}[1]{>{\centering\arraybackslash}p{#1}}

% ============================================================
\begin{document}
\pagenumbering{gobble}   % suppress page number on title page

% ── TITLE PAGE ───────────────────────────────────────────────
\begin{titlepage}
    \centering

    \begin{figure}[H]
        \centering
        \includegraphics[width=5cm]{ieu-logo.png}
        \label{fig:logo}
    \end{figure}

    \vspace{0.6cm}
    {\large\bfseries T.C.}\\[0.3cm]
    {\large\bfseries IZMIR UNIVERSITY OF ECONOMICS}\\[0.2cm]
    {\large FACULTY OF ENGINEERING}\\[0.8cm]

    \rule{\linewidth}{0.6pt}\\[0.4cm]
    {\LARGE\bfseries CE 216}\\[0.3cm]
    {\Large\bfseries ``SPORTS MANAGER'' PROJECT}\\[0.3cm]
    {\Large\bfseries TEAM: STACK UNDERFLOW}\\[0.3cm]
    {\large\bfseries DESIGN AND ARCHITECTURE DOCUMENT}\\[0.2cm]
    \rule{\linewidth}{0.6pt}\\[1.2cm]

    \begin{tabular}{ll}
        \textbf{S\.{I}MGE GÖKALP}        & 20230602029 \\[0.3cm]
        \textbf{TUANA YA\u{G}MUR}         & 20240602410 \\[0.3cm]
        \textbf{V\.{I}LDAN TURNALAR \.{I}NAL} & 20253501001 \\[0.3cm]
        \textbf{U\u{G}UR EM\.{I}N BAYNAL} & 20220602408 \\
    \end{tabular}

    \vfill
    {\large\textbf{LECTURER:} KAYA O\u{G}UZ}\\[0.5cm]
    {\large \the\year}
\end{titlepage}

% ── TABLE OF CONTENTS ────────────────────────────────────────
\pagenumbering{roman}
\tableofcontents
\newpage
\pagenumbering{arabic}

% ============================================================
%  1. INTRODUCTION
% ============================================================
\section{Introduction}

Sports Manager application is a GUI-based management application that allows a user to
manage a team within a league-based sports competition.

The application is designed as a generic sports management framework that supports multiple
sports. The system architecture allows new sports to be added without modifying the core
application or the user interface, ensuring extensibility through Object-Oriented Design
principles.

Initially, Football will be implemented as the first sport, while the architecture supports the
future addition of other sports.

% ============================================================
%  2. SCOPE AND FUNCTIONALITY
% ============================================================
\section{Scope and Functionality}

The system simulates a sports league and allows the user to manage a team within that
league. The project includes the following main components:

\subsubsection*{League Management}
\begin{itemize}
    \item Generation of a league consisting of multiple teams
    \item Creation of players and coaches with randomized attributes
    \item Generation of a full season fixture schedule and weekly progression
    \item Maintenance of league standings
\end{itemize}

\subsubsection*{Team Management}
The user will act as a team manager and will be able to:
\begin{itemize}
    \item View team roster
    \item Select starting lineup
    \item Assign substitutes
    \item Configure team tactics
    \item Monitor player injuries and availability
\end{itemize}

\subsubsection*{Match Simulation}
Matches will be simulated automatically by the system.
Key characteristics:
\begin{itemize}
    \item Matches are processed in segments (halves, quarters, or sets depending on the sport)
    \item The user cannot directly control the players
    \item Tactical adjustments and substitutions can be made between segments
\end{itemize}

\subsubsection*{Injury Management}
Players may become injured during matches. Injuries are tracked based on the number of
matches missed, not time duration.

% ─────────────────────────────────────────────────────────────
\subsection{Functional Requirements}

The functional requirements are listed below in user story format and are organized by system
components, which are reflected in the requirement ID prefixes.

\subsubsection*{League Management}

\begin{longtable}{L{1.4cm} L{2.8cm} L{8.5cm}}
\toprule
\textbf{ID} & \textbf{Functional Requirement} & \textbf{User Story} \\
\midrule
\endhead
LM 1 & Start a New Game &
    As a player, I want to start a new game so that I can begin managing a sports team in a new season. \\[0.4em]
LM 2 & Select a Sport &
    As a player, I want to choose the sport at the beginning of the game so that I can play the manager simulation for different sports. \\[0.4em]
LM 3 & Generate a League &
    As a player, I want the system to generate a league with multiple teams so that I can play in a complete season competition. \\[0.4em]
LM 4 & Generate Teams, Players, and Coaches &
    As a player, I want the system to generate teams with players and coaches so that each new game starts with a complete playable league. \\[0.4em]
LM 5 & Support Sport-Specific Attributes and Positions &
    As a player, I want players to have sport-specific positions and attributes so that each sport feels different and realistic. \\[0.4em]
LM 6 & Select a Team to Manage &
    As a player, I want to choose one team to manage so that I can control its lineup, tactics, and season progress. This is implied by the manager game flow described in the project brief. \\[0.4em]
LM 7 & View the League Table &
    As a player, I want to see the league standings so that I can track my team's position during the season. \\[0.4em]
LM 8 & Update League Standings After Matches &
    As a player, I want the league table to be updated after each match so that I can see the current rankings. \\[0.4em]
LM 9 & Finish a Season &
    As a player, I want the system to determine whether I won the league at the end of the season so that I can see the outcome of my performance. \\[0.4em]
LM 10 & Load Names and Visuals &
    As a player, I want team names, player names, and visual logos to be loaded from resources so that the game has a good presentation. \\[0.4em]
LM 11 & Add New Sports &
    As a developer, I want to add a new sport without rewriting the existing UI so that the framework remains extensible. \\
\bottomrule
\caption{League Management Functional Requirements}
\end{longtable}

\subsubsection*{Team Management}

\begin{longtable}{L{1.4cm} L{2.8cm} L{8.5cm}}
\toprule
\textbf{ID} & \textbf{Functional Requirement} & \textbf{User Story} \\
\midrule
\endhead
TM 1 & View the Player List &
    As a player, I want to see the list of players in my team so that I can make squad and tactical decisions. \\[0.4em]
TM 2 & View the Team Schedule &
    As a player, I want to view my team's fixture schedule so that I can follow upcoming and past matches. \\[0.4em]
TM 3 & Set Lineup Before a Match &
    As a player, I want to choose the team lineup before a match so that I can decide who will play. \\[0.4em]
TM 4 & Set Tactics Before a Match &
    As a player, I want to set tactics before a match so that I can influence how my team performs. \\
\bottomrule
\caption{Team Management Functional Requirements}
\end{longtable}

\subsubsection*{Match Simulation}

\begin{longtable}{L{1.4cm} L{2.8cm} L{8.5cm}}
\toprule
\textbf{ID} & \textbf{Functional Requirement} & \textbf{User Story} \\
\midrule
\endhead
MS 1 & Simulate Matches Automatically &
    As a player, I want matches to be simulated by the system so that I can focus on managerial decisions rather than directly controlling players. \\[0.4em]
MS 2 & Play Matches in Segments &
    As a player, I want matches to progress in segments such as quarters, halves, or sets so that I can intervene at break points depending on the sport. \\[0.4em]
MS 3 & Change Tactics During Breaks &
    As a player, I want to change tactics between match segments so that I can respond to the state of the match. \\[0.4em]
MS 4 & Substitute Players During Breaks &
    As a player, I want to substitute players during break points so that I can react to injuries, fatigue, or tactics. \\[0.4em]
MS 5 & Apply Sport-Specific Scoring Rules &
    As a player, I want each sport to apply its own scoring and points rules so that league results are calculated correctly. \\
\bottomrule
\caption{Match Simulation Functional Requirements}
\end{longtable}

\subsubsection*{Injury Management}

\begin{longtable}{L{1.4cm} L{2.8cm} L{8.5cm}}
\toprule
\textbf{ID} & \textbf{Functional Requirement} & \textbf{User Story} \\
\midrule
\endhead
IM 1 & Track Injuries by Number of Games &
    As a player, I want injured players to miss a specific number of matches so that injury management affects squad selection. \\[0.4em]
IM 2 & Prevent Injured Players from Being Selected &
    As a player, I want injured players to be unavailable for selection so that invalid match lineups cannot be created. \\
\bottomrule
\caption{Injury Management Functional Requirements}
\end{longtable}

% ─────────────────────────────────────────────────────────────
\subsection{Non-Functional Requirements}

The non-functional requirements are listed below and describe the quality attributes and
architectural constraints of the system.

\begin{description}[leftmargin=2em, labelwidth=2em]
    \item[\textbf{NFR 01} -- Graphical User Interface:]
        The system shall provide a graphical user interface that allows users to interact with the sports management simulation easily.
    \item[\textbf{NFR 02} -- Extensibility:]
        The system shall support the addition of new sports without requiring modifications to the existing user interface or core logic.
    \item[\textbf{NFR 03} -- Maintainability:]
        The system shall be structured using interfaces, abstract classes, and modular components to ensure maintainable and readable code.
    \item[\textbf{NFR 04} -- Testability:]
        Core framework components shall be designed in a way that allows unit testing of their functionality.
    \item[\textbf{NFR 05} -- Reusability:]
        Common domain entities such as Team, Player, League, and Match shall be reusable across different sports.
    \item[\textbf{NFR 06} -- Installability:]
        The final application shall be installable on Windows systems.
\end{description}

% ============================================================
%  3. SYSTEM ARCHITECTURE PRINCIPLES
% ============================================================
\section{System Architecture Principles}

The system architecture of the Sports Manager Framework is designed according to key
Object-Oriented Design principles in order to ensure extensibility and maintainability.
Since the application must support multiple sports without modifying the core system or the
user interface, the architecture separates generic game logic from sport-specific behavior.
The following design principles and patterns are planned to be used in the system to
implement such a system.

\subsection{Interface-Based Design}

The system architecture follows an interface-based design approach, where the main
components interact through abstract interfaces. In this approach, the user interface
communicates with the core system through generic abstractions such as Sport, League, and
Player. As a result, the user interface does not depend on a specific sport implementation.
Instead, it interacts with the system through common interfaces.

\subsection{Factory Pattern}

The Factory Pattern is used to create sport-specific objects at runtime without exposing the
object creation logic to the rest of the system.

When the user starts a new game and selects a sport, the application uses a factory
component to create the appropriate sport module:
\begin{itemize}
    \item[\textrightarrow] \texttt{FootballSport}
\end{itemize}

Instead of directly instantiating class in the UI, the system calls a factory method such as:
\begin{lstlisting}
Sport sport = SportFactory.createSport(selectedSport);
\end{lstlisting}

This design keeps the application independent from specific sport implementations.

\subsection{Layered Architecture}

The architecture separates the system into two main layers as Presentation Layer (User
Interface) and Core Domain Layer which is explained in the System Architecture section.
This separation ensures that the user interface remains independent from sport-specific logic.

% ============================================================
%  4. SYSTEM ARCHITECTURE
% ============================================================
\section{System Architecture}

\subsection{Overview}

The Sports Manager application is designed based on a multi-layered structure that provides a
clear separation between the UI (user interface), the core game logic, and sport-specific
implementation. This system mainly aims to design a UI layer that is independent of any
sport-specific classes. To achieve this, the communication between layers is grounded on
abstractions. The abstraction ensures that when a new sport is added to the Sports Manager
system, the existing UI design will not be affected, and it will be applicable for different sports.

In this project, the architecture has four distinct layers: User Interface, Core System, Sport
Abstraction, and Concrete Sport Module. The high-level layer structure is illustrated in
Figure~\ref{fig:highlevel}, and the detailed class relationships are shown in Figure~\ref{fig:classdesign}
(Section~8.1). Between these layers, there is one-directional downward dependency, which
means that upper classes are not allowed to directly import or instantiate any classes from
lower layers. The only way to communicate with upper layers is via the interfaces and
abstractions defined in the Sport Abstraction layer. This limitation is the main logic of the
system, which makes it extensible by design rather than by coincidence.

You can see the high level design below represented in a class diagram format.

\begin{figure}[H]
    \centering
    \includegraphics[width=\textwidth]{figure1-highlevel.png}
    \caption{High-Level Architecture Layer Structure}
    \label{fig:highlevel}
\end{figure}

\subsection{Dependency Rule}

Before describing each layer in detail, it is crucial to understand the fundamental architectural
principle of the system:

\begin{itemize}
    \item The UI layer only depends on the SportManager
    \item SportsManager depends on GameSession, League, Team, and the Sport interface.
    \item The Core System layer depends exclusively on the Sport interface and never on any concrete sport implementation such as FootballSport
    \item Concrete Sport Modules depend on the Sport Abstraction layer by implementing and extending its interfaces and abstract classes.
\end{itemize}

This means that when a developer adds a new sport, they only write new classes in the
Concrete Sport Module. The rest of the system never needs to be changed.

\subsection{User Interface Layer}

The user interface layer includes JavaFX controllers that operate screen rendering and
interactions with the user. The controllers in this layer are:

\begin{itemize}
    \item \textbf{MainMenuController:} It displays the main menu and manages starting a new game or loading the previous one.
    \item \textbf{SportSelectionController:} It is shown at the beginning of a new game and lists available sports and sends the user's chosen sport code to SportManager.
    \item \textbf{DashboardController:} It is the main game screen, which shows weekly progression and calls \texttt{SportManager.advanceWeek()} to move the game forward.
    \item \textbf{SquadController:} It retrieves the player list from the management team and shows the team roster, giving the user access to player attributes and availability.
    \item \textbf{MatchController:} It controls pre-match setup (lineup, tactics) and displays match segments and outcomes.
    \item \textbf{StandingsController:} It uses \texttt{SportManager.showLeagueTable()} to retrieve and display the current league table.
    \item \textbf{FixtureController:} It retrieves and displays the full season fixture schedule by calling \texttt{SportManager.showFixture()}.
\end{itemize}

Importantly, this layer does not contain game logic. All operations are delegated exclusively
through SportManager. The UI controllers never reference sport-specific classes. They only
interact with the Sport interface from the Sport Abstraction layer and the generic types Team,
Player, Match, League, and MatchResult from the Core System layer. This implies that
regardless of the sport being played, the user interface functions exactly the same.

\subsection{Core System Layer}

The Core System layer contains the main game management logic and all sport-independent
domain entities. It is responsible for running the game loop, managing the season, and
organizing all game components. The classes in this layer do not reference any concrete sport
implementation.

\subsubsection*{a) SportManager}

SportManager is the main entry point of the game application and the only class that
the UI layer directly interacts with. It serves as a front, masking the complexity of the
remainder of the system behind an easy-to-use interface. The following are the key
methods of SportManager:

\begin{itemize}
    \item \texttt{startNewGame():} Initializes a new GameSession and triggers sport selection.
    \item \texttt{selectSport(sportCode: String):} Calls \texttt{SportFactory.createSport(sportCode)} to obtain the correct Sport implementation and stores it in the session.
    \item \texttt{selectManagedTeam(teamId: String):} Sets the team the user will manage for the season.
    \item \texttt{advanceWeek(week: String):} Progresses the game by one week; triggers training and schedules upcoming matches.
    \item \texttt{startMatch(matchId: String):} Retrieves the match for the given ID, delegates simulation to the sport-specific Match class, and returns the result to the UI.
    \item \texttt{showLeagueTable():} Retrieves the current standings from the League object.
    \item \texttt{showFixture():} Retrieves the full fixture list from the League object.
\end{itemize}

\subsubsection*{b) GameSession}

GameSession holds the complete state of the currently running game. It stores:

\begin{itemize}
    \item \texttt{currentSeasonYear:} The current season number
    \item \texttt{currentWeek:} The current week within the season
    \item \texttt{selectedSport:} The active Sport implementation (stored as the interface type, never as a concrete class)
    \item \texttt{managedTeam:} The team the user is managing
    \item \texttt{league:} The active League object for the current season
\end{itemize}

It also provides \texttt{isSeasonFinished()} to check whether all fixtures have been played and
\texttt{startNextSeason()} to reset the state for a new season while keeping the same sport
and managed team.

\subsubsection*{c) League}

League manages the full season structure. It is responsible for:

\begin{itemize}
    \item Storing the list of all teams in the season
    \item Generating the full fixture schedule through \texttt{generateFixtures()}, which creates a round-robin schedule where each team plays every other team twice
    \item Retrieving all matches for a given week through \texttt{getMatchesOfWeek(weekNo: int)}, which returns a \texttt{List<Match>}
    \item Updating the standings after a match through \texttt{updateStandings(result: MatchResult)}, which adjusts the relevant StandingEntry objects
    \item Exposing the current league table through \texttt{getTable()}, which returns a sorted \texttt{List<StandingEntry>}
\end{itemize}

\subsubsection*{d) Team}

Team represents a sports team and holds the following state:

\begin{itemize}
    \item \texttt{name} and \texttt{logoPath} for display purposes
    \item \texttt{players: List<Player>}: The full squad
    \item \texttt{coaches: List<Coach>}: The coaching staff
    \item \texttt{currentLineup: Lineup}: The currently selected lineup
    \item \texttt{currentTactic: Tactic}: The currently selected tactic
\end{itemize}

\subsubsection*{e) Match}

Match is an abstract class representing a match between two teams. It holds:

\begin{itemize}
    \item \texttt{weekNo}, \texttt{homeTeam}, \texttt{awayTeam}: Match context
    \item \texttt{homeScore}, \texttt{awayScore}: running scores updated during simulation
    \item \texttt{segments: List<MatchSegment>}: The list of match segments (halves, quarters, or sets depending on the sport)
    \item \texttt{result: MatchResult}: The final outcome
\end{itemize}

It defines three abstract methods that sport-specific subclasses must implement:

\begin{itemize}
    \item \texttt{start():} Initializes the match and prepares segment structure
    \item \texttt{play():} Simulates the next segment and updates scores and events
    \item \texttt{finish():} Finalizes the result, determines the winner, and applies injuries
\end{itemize}

\subsubsection*{f) Player}

Player is an abstract class representing a player. It holds:

\begin{itemize}
    \item \texttt{name}, \texttt{age}, \texttt{position}: Identification and role
    \item \texttt{fitness}: Current physical condition, affects match performance
    \item \texttt{injuredMatches}: Number of matches the player must miss due to injury
\end{itemize}

It defines three abstract methods:

\begin{itemize}
    \item \texttt{isAvailable(): boolean}: Returns whether the player can be selected for a match
    \item \texttt{train():} Applies a training session, adjusting fitness and sport-specific attributes
    \item \texttt{recover():} Decrements the injury counter and restores availability when the counter reaches zero
\end{itemize}

\subsubsection*{Supporting Classes}

\begin{itemize}
    \item \textbf{Coach:} holds coaching staff information: name, role, shape, training skill, and motivation skill. Coaches affect the outcome of training sessions.
    \item \textbf{Lineup:} holds the list of starters and substitutes for a match. It provides \texttt{isValid(): boolean}, which checks whether the lineup satisfies the sport-specific constraints such as minimum and maximum player counts.
    \item \textbf{Tactic:} holds the tactic configuration: name, shape, offense level, and defense level.
    \item \textbf{MatchResult:} stores the final score, the winning team reference, and the list of InjuryRecord objects produced during the match.
    \item \textbf{MatchSegment:} represents a single segment of a match. It holds the segment number, a display label such as ``First Half'' or ``Quarter 3'', partial scores for home and away, and a list of event strings describing what happened during that segment.
    \item \textbf{InjuryRecord:} tracks a player injury. It stores the player reference, the number of games remaining, a description of the injury, and provides \texttt{isRecovered(): boolean} to check whether the player is ready to return.
    \item \textbf{StandingEntry:} tracks a single team's league record, storing played, won, drawn, lost, and points values for standings calculation and display.
\end{itemize}

\subsection{Sport Abstraction Layer}

The agreement that each sport must uphold is outlined in the Sport Abstraction layer. It is the
line separating any sport-specific implementation from the general core system. This layer is
necessary for the Core System layer. This layer is implemented by the Concrete Sport
Module.

\subsubsection*{a) Sport Interface}

Sport is a Java interface that declares all sport-specific behaviors. The complete set of
methods is:

\begin{itemize}
    \item \texttt{getName(): String}: Returns the display name of the sport, used in the UI
    \item \texttt{createLeague(teamCount: int): League}: Generates and returns a fully populated league with the given number of teams, including all teams, players, and coaches
    \item \texttt{createTeam(name: String): Team}: Creates and returns a sport-specific team instance
    \item \texttt{createPlayer(name: String, position: String): Player}: Creates and returns a sport-specific player with appropriate attributes for the given position
    \item \texttt{createMatch(home: Team, away: Team, weekNo: int): Match}: Creates and returns a sport-specific match instance
    \item \texttt{calculatePoints(result: MatchResult, team: Team): int}: Calculates and returns the points earned by the given team from the given result, applying sport-specific scoring rules
    \item \texttt{trainTeam(team: Team):} Applies a full training session to the team, updating player fitness and attributes
    \item \texttt{applyInjuries(match: Match):} Processes the match result and generates injury records for affected players
\end{itemize}

Every method in this interface is called by the Core System layer during normal game
operation. Because the Core System only ever holds a reference of type Sport, it works
identically for any sport that implements this interface.

\subsubsection*{b) SportFactory}

SportFactory is responsible for instantiating the correct sport implementation at
runtime. It exposes a single static method:

\begin{itemize}
    \item \texttt{createSport(sportCode: String): Sport}
\end{itemize}

When the user selects a sport at the start of a new game, SportManager calls this
method with the selected sport code. The factory returns the appropriate Sport
implementation. Because the return type is the Sport interface, the caller never knows
or needs to know which concrete class was created. Adding a new sport to the factory
requires adding a single new case to this method.

\subsection{Concrete Sport Module}

The Concrete Sport Module contains all sport-specific implementations. For the current
version of the application, this module contains the Football implementation.

\begin{itemize}
    \item \textbf{FootballSport:} Implements the Sport interface and defines all Football-specific rules and configurations. It determines how Football leagues are structured, how Football players are generated with Football-specific positions such as goalkeeper, defender, midfielder, and forward, how Football matches are simulated in two halves, and how points are awarded — three for a win, one for a draw, and zero for a loss.

    \item \textbf{FootballLeague:} Extends League and adds Football-specific standing tie-breaking rules. When two teams have equal points, it compares head-to-head results, then goal difference, then a coin toss, exactly as specified in the project requirements.

    \item \textbf{FootballMatch:} Extends Match and implements \texttt{start()}, \texttt{play()}, and \texttt{finish()} for a two-half Football match. It handles Football-specific scoring, event generation, and injury application.

    \item \textbf{FootballTeam:} Extends Team and adds Football-specific team behavior including formation validation and Football-specific tactic constraints.

    \item \textbf{FootballPlayer:} Extends Player and adds Football-specific attributes such as shooting, passing, defending, and speed. It implements \texttt{isAvailable()}, \texttt{train()}, and \texttt{recover()} with Football-specific logic.
\end{itemize}

\subsection{Data Flow: Starting a Game and Playing a Match}

To make the architecture concrete, the following describes the complete data flow for two key
scenarios.

\subsubsection*{Starting a new game:}

\begin{enumerate}
    \item The user clicks ``New Game'' in MainMenuController, which calls \texttt{SportManager.startNewGame()}
    \item SportManager creates a new GameSession and navigates to SportSelectionController
    \item The user selects Football; SportSelectionController calls \texttt{SportManager.selectSport("football")}
    \item SportManager calls \texttt{SportFactory.createSport("football")}, which returns a FootballSport instance stored as type Sport in GameSession
    \item SportManager calls \texttt{sport.createLeague(20)}, which generates 20 Football teams with players and coaches
    \item The league is stored in GameSession and DashboardController is displayed
\end{enumerate}

\subsubsection*{Playing a match:}

\begin{enumerate}
    \item The user selects a match from FixtureController and confirms the lineup in MatchController
    \item MatchController calls \texttt{SportManager.startMatch(matchId)}
    \item SportManager retrieves the match from \texttt{League.getMatchesOfWeek()} and calls \texttt{match.start()}
    \item For each segment, \texttt{match.play()} is called, generating events and updating scores
    \item \texttt{match.finish()} is called, producing a MatchResult and triggering \texttt{sport.applyInjuries(match)}
    \item SportManager calls \texttt{league.updateStandings(result)} and returns the result to MatchController
    \item MatchController displays the match segments, final score, and any injuries to the user
\end{enumerate}

% ============================================================
%  5. CORE INTERFACES
% ============================================================
\section{Core Interfaces}

This section will cover the interface structures that the application uses. The main reason for
using these interface structures is that the system can run without being based on a specific
sport. In other words, the application is independent of football and other sports. Instead, the
general characteristics of different sports are defined in these interface structures. The classes
that are created for different sports will be based on these interface structures and will include
their specific characteristics and rules. The main advantage of this approach is that the system
can be easily extended. If a new sport has to be included in the system in the future, there is
no need to make changes to a large part of the code. Instead, new classes can be created
that will be based on these interface structures. This approach is also in agreement with the
Open/Closed Principle. In other words, the system is closed to changes and open to
extensions.

\subsection{Sport Interface}

The Sport interface allows different sports to be represented in our system. Even though each
sport has different rules and characteristics, they share common operations.

These operations are defined within the Sport interface. These operations are performed
according to different rules for each sport. For example, our system has a FootballSport class
for football. The Sport interface basically defines the following operations:

\begin{itemize}
    \item Getting the name of the sport
    \item Creating a league
    \item Creating a team
    \item Creating a match
    \item Updating points based on the match result
\end{itemize}

This structure allows the system to support different sports, and other parts of the application
can function without being tied to the specifics of a particular sport.

As described in the System Architecture Principles section, Factory pattern will be used to
create sport specific classes.
\begin{figure}[H]
    \centering
    \includegraphics[width=\textwidth]{sport abstraction}
    \caption{Sport Abstraction}
    \label{fig:abstractclasses}
\end{figure}

% ============================================================
%  6. ABSTRACT CLASSES
% ============================================================
\section{Abstract Classes}

Abstract classes are used to represent the basic structures that are common to many sports.
The abstract classes define the features that are common to many sports.

With the use of abstract classes, the code can be written only once, as opposed to repeating
the code. The specific classes for the respective sports can extend the classes to make the
required modifications.

The basic abstract classes used for the project are: Player, Team, Match and League.
These classes represent the core assets of the sports management system.

The class diagram for these classes are represented below.

\begin{figure}[H]
    \centering
    \includegraphics[width=\textwidth]{figure2-abstract-classes.png}
    \caption{Abstract Classes Class Diagram (Player, Team, Match, League)}
    \label{fig:abstractclasses}
\end{figure}

\subsection{Player Abstract Class}

The Player class describes the overall structure of the players in the system. Information such
as a player's name, position, skill level, and injury status is maintained in this class.

The FootballPlayer class, designed for the sport of football, can be a subclass of this class and
can add additional data.

\subsection{Team Abstract Class}

The overall structure of a team can be defined by a class named Team. Players, coaching
staff, and tactical details can be included in this class.

The FootballTeam class can be a derived class that includes specific details for football teams.

\subsection{Match Abstract Class}

The Match class represents a game played between two teams. Information about which
teams played the game and the score is stored in this class.

The FootballMatch class, created for football, inherits from this class and implements rules
specific to football matches.

\subsection{League Abstract Class}

The League abstract class models the core structure of a sports competition, including teams,
fixtures, and league standings. It defines common league operations such as fixture
generation, match scheduling, and standings updates, while allowing sport-specific
implementations to extend its behavior.

% ============================================================
%  7. SPORT IMPLEMENTATION (FOOTBALL EXAMPLE)
% ============================================================
\section{Sport Implementation (Football Example)}

In addition, the design of the system includes abstraction and inheritance concepts for
different sports. This section will discuss how the system handles the sport of football as a
case study.

In the system, the sport of football is implemented in the system by using the FootballSport
classes, which are derived from the interface class called Sport. Basic operations concerning
the sport of football can be performed in these classes.

In addition, other classes concerning the sport of football in the system are derived from the
above abstract classes using the concept of inheritance. For example, the FootballPlayer
class is derived from the abstract class called Player, which contains properties concerning
football players. In addition, the FootballTeam class, derived from the abstract class called
Team, contains properties concerning football teams. Furthermore, the FootballMatch class,
derived from the abstract class called Match, contains properties concerning football matches.
Finally, the FootballLeague class, derived from the abstract class called League, contains
properties concerning football leagues.

In this way, the sport of football can be easily incorporated into the system.
\begin{figure}[H]
    \centering
    \includegraphics[width=0.65\textwidth]{figure3-core-layer.png}
    \caption{Sport Implementation UML}
    \label{fig:corelayer}
\end{figure}

% ============================================================
%  8. DETAILED DESIGN: CLASSES AND METHODS
% ============================================================
\section{Detailed Design: Classes and Methods}

\subsection{Core Layer}

\begin{wrapfigure}{r}{0.60\textwidth}
    \centering
    \includegraphics[width=0.58\textwidth]{figure3-core-layer1.png}
    \caption{Core Layer — SportManager, GameSession, and Domain Entities}
    \label{fig:corelayer}
\end{wrapfigure}

The Core module represents the central part of the Sports Manager framework and contains
the main domain model and application logic. It includes key entities which model the
structure of a sports competition. SportManager and GameSession classes coordinate the
overall game flow, including league progression, match execution, and team management.
This separation ensures that the core logic remains independent from the user interface and
sport-specific implementations.

\clearpage
\subsection{Overall Package and Class Design}

The diagram below shows the full picture of classes and relations between main modules.





\begin{figure}[H]
    \centering
    \includegraphics[width=\textwidth]{figure4-full-class-design.png}
    \caption{Full Package and Class Relationships}
    \label{fig:classdesign}
\end{figure}

\subsection{Draft UI Classes}

\begin{wrapfigure}{r}{0.60\textwidth}
    \centering
    \includegraphics[width=0.58\textwidth]{UIlayer.png}
    \caption{UI Layer}
    \label{fig:corelayer}
\end{wrapfigure}

UI controllers will be aligned with the graphical UI designs during the implementation phase.
\clearpage




% ============================================================
%  9. DETAILED DESIGN: TRACEABILITY TO REQUIREMENTS
% ============================================================
\section{Detailed Design: Traceability to Requirements}

The requirements are mapped to the corresponding classes and methods to ensure
traceability between the system requirements and the detailed design. The mappings are
presented in the tables below.

\subsubsection*{League Management Mapping}

{\footnotesize
\begin{longtable}{L{0.9cm} L{2.0cm} L{2.1cm} L{5.0cm} L{3.5cm}}
\toprule
\textbf{ID} & \textbf{Functional Req.} & \textbf{Primary Class} & \textbf{Related Method(s)} & \textbf{Explanation} \\
\midrule
\endhead
LM-1 & Start a New Game        & SportManager            & \texttt{startNewGame()}                          & Starts a new game session and initializes the simulation state. \\[0.4em]
LM-2 & Select a Sport          & SportManager, SportFactory & \texttt{selectSport()}\newline\texttt{createSport()}   & Creates the selected sport implementation through the factory. \\[0.4em]
LM-3 & Generate a League       & Sport, League           & \texttt{createLeague()}\newline\texttt{addTeam()}\newline\texttt{generateFixtures()} & Creates the league structure and generates the fixture schedule. \\[0.4em]
LM-4 & Generate Teams, Players, and Coaches & Sport, Team & \texttt{createTeam()}\newline\texttt{createPlayer()}\newline\texttt{addPlayer()} & Creates teams and populates them with players and coaches. \\[0.4em]
LM-5 & Support Sport-Specific Attributes & Sport, Player & \texttt{createPlayer(name, position)} & Creates sport-specific player instances with appropriate attributes. \\[0.4em]
LM-6 & Select a Team to Manage & SportManager            & \texttt{selectManagedTeam(teamId)} & Assigns the selected team to the active game session. \\[0.4em]
LM-7 & View the League Table   & SportManager, League    & \texttt{showLeagueTable()}\newline\texttt{getTable()}  & Displays the current league standings. \\[0.4em]
LM-8 & Update League Standings & League, Sport           & \texttt{updateStandings(result)}\newline\texttt{calculatePoints()} & Updates standings and points after each match. \\[0.4em]
LM-9 & Finish a Season         & League, GameSession     & \texttt{advanceWeek()}\newline\texttt{getMatchesOfWeek()} & Season completion is inferred when all fixtures are played. \\[0.4em]
LM-10 & Load Names and Visuals & —                       & —                           & Not represented directly in the UML; can be implemented with a ResourceLoader. \\[0.4em]
LM-11 & Add New Sports         & SportFactory            & \texttt{createSport(sportCode)} & Allows the addition of new sport modules without modifying the UI. \\
\bottomrule
\caption{League Management Traceability}
\end{longtable}
}

\subsubsection*{Team Management Mapping}

{\footnotesize
\begin{longtable}{L{0.9cm} L{2.0cm} L{2.1cm} L{5.0cm} L{3.5cm}}
\toprule
\textbf{ID} & \textbf{Functional Req.} & \textbf{Primary Class} & \textbf{Related Method(s)} & \textbf{Explanation} \\
\midrule
\endhead
TM-1 & View the Player List   & Team                       & \texttt{getAvailablePlayers()} & Returns the list of players currently available for selection. \\[0.4em]
TM-2 & View the Team Schedule & SportManager, League       & \texttt{showFixture()}\newline\texttt{getMatchesOfWeek()} & Displays the team's upcoming and past matches. \\[0.4em]
TM-3 & Set Lineup Before a Match & Team, Lineup            & \texttt{setLineup()}\newline\texttt{isValid()} & Defines and validates the starting lineup and substitutes. \\[0.4em]
TM-4 & Set Tactics Before a Match & Team                   & \texttt{setTactic(tactic)} & Allows the manager to configure tactical parameters. \\
\bottomrule
\caption{Team Management Traceability}
\end{longtable}
}

\subsubsection*{Match Simulation Mapping}

{\footnotesize
\begin{longtable}{L{0.9cm} L{2.0cm} L{2.1cm} L{5.0cm} L{3.5cm}}
\toprule
\textbf{ID} & \textbf{Functional Req.} & \textbf{Primary Class} & \textbf{Related Method(s)} & \textbf{Explanation} \\
\midrule
\endhead
MS-1 & Simulate Matches Automatically & SportManager, Match & \texttt{startMatch()}\newline\texttt{start()}\newline\texttt{play()}\newline\texttt{finish()} & Runs the automated match simulation. \\[0.4em]
MS-2 & Play Matches in Segments & Match, MatchSegment     & \texttt{play()}            & Processes the match flow through multiple segments. \\[0.4em]
MS-3 & Change Tactics During Breaks & Team                 & \texttt{setTactic()}       & Allows tactical adjustments during match breaks. \\[0.4em]
MS-4 & Substitute Players During Breaks & Team, Lineup     & \texttt{setLineup()}       & Allows substitutions during match intervals. \\[0.4em]
MS-5 & Apply Sport-Specific Scoring Rules & Sport          & \texttt{calculatePoints()} & Applies sport-specific scoring logic. \\
\bottomrule
\caption{Match Simulation Traceability}
\end{longtable}
}

\subsubsection*{Injury Management Mapping}

{\footnotesize
\begin{longtable}{L{0.9cm} L{2.0cm} L{2.1cm} L{5.0cm} L{3.5cm}}
\toprule
\textbf{ID} & \textbf{Functional Req.} & \textbf{Primary Class} & \textbf{Related Method(s)} & \textbf{Explanation} \\
\midrule
\endhead
IM-1 & Track Injuries by Number of Games & Sport, InjuryRecord, Player & \texttt{applyInjuries()}\newline\texttt{decrementGames()}\newline\texttt{recover()} & Tracks injuries based on the number of matches missed. \\[0.4em]
IM-2 & Prevent Injured Players from Being Selected & Player, Team & \texttt{isAvailable()}\newline\texttt{getAvailablePlayers()} & Ensures injured players cannot be selected for matches. \\
\bottomrule
\caption{Injury Management Traceability}
\end{longtable}
}

% ============================================================
%  10. DETAILED DESIGN: USER INTERFACES
% ============================================================
\section{Detailed Design: User Interfaces}

The Sports Manager application provides a graphical user interface built with JavaFX,
following a dark sports-management aesthetic with a deep navy palette and amber accent
colours. All screens share a consistent visual language: a top navigation bar, card-based
content panels, and a unified CSS theme applied globally through \texttt{style.css}. The UI is
strictly read-only with respect to the domain — every screen only calls methods on
SportManager and renders what is returned through abstract types. No controller is aware of
which concrete sport is loaded.

\subsection{Main Menu Screen}

\textit{Controller:} \texttt{MainMenuController} \quad|\quad \textit{FXML:} \texttt{main-menu.fxml}

The Main Menu is the application's entry point. It is displayed immediately when the
application launches. The screen features a full-height centred layout with a large football
emoji icon, the application title ``SPORT MANAGER'', and a subtitle line. Below these, three
vertically stacked action buttons are presented:

\begin{itemize}
    \item \textbf{New Game} — starts a new game session by calling \texttt{SportManager.startNewGame()}, which resets the GameSession and navigates to the Sport Selection screen.
    \item \textbf{Load Game} — reserved for future functionality; currently displays an informational dialog.
    \item \textbf{Exit} — closes the application.
\end{itemize}

The design purpose of this screen is minimal: it provides a clean, branded starting point and
delegates all responsibility immediately to SportManager.

\begin{figure}[H]
    \centering
    \includegraphics[width=0.8\textwidth]{screen-main-menu.png}
    \caption{Main Menu Screen}
    \label{fig:screen-mainmenu}
\end{figure}

\subsection{Sport Selection Screen}

\textit{Controller:} \texttt{SportSelectionController} \quad|\quad \textit{FXML:} \texttt{sport-selection.fxml}

The Sport Selection screen is shown after the user starts a new game. Its purpose is to allow
the user to choose which sport they wish to manage for the season (Requirement LM-2). The
screen header displays the title ``SELECT YOUR SPORT''. The body contains a scrollable
FlowPane populated dynamically by the controller at runtime — one card is generated per
available sport.

Each sport card contains:

\begin{itemize}
    \item A large sport emoji icon
    \item The sport's display name in bold
    \item A short description line beneath it
\end{itemize}

Currently, Football is the only enabled sport, styled with full opacity and a hover glow effect.
Future sports (Basketball, Tennis) are displayed at reduced opacity with the cursor blocked,
visually communicating that they are planned but not yet implemented. When the user clicks
an enabled card, \texttt{SportManager.selectSport(code)} is called, which triggers league generation
and navigates forward.

A Back button in the bottom bar returns the user to the Main Menu.

\begin{figure}[H]
    \centering
    \includegraphics[width=0.8\textwidth]{screen-sport-selection.png}
    \caption{Sport Selection Screen}
    \label{fig:screen-sportselection}
\end{figure}

\subsection{Team Selection Screen}

\textit{Controller:} \texttt{TeamSelectionController} \quad|\quad \textit{FXML:} \texttt{team-selection.fxml}

The Team Selection screen allows the user to choose which team they will manage for the
season (Requirement LM-6). The top bar displays the generated league name and a subtitle
showing how many clubs are available.

The centre of the screen contains a scrollable GridPane arranged in four columns. Each cell
contains a team card displaying:

\begin{itemize}
    \item The team's name in bold
    \item The number of players in the roster and their average skill score
    \item The number of coaching staff members
\end{itemize}

Teams are generated by \texttt{FootballSport.createLeague(20)} before this screen is shown, so the
user sees all 20 randomly named clubs with realistic roster statistics. Clicking a card selects it,
which highlights its border in amber. Clicking a different card deselects the previous one. The
user confirms their choice with the Confirm Selection button at the bottom, which calls
\texttt{SportManager.selectManagedTeam(team)} and navigates to the Dashboard.

\begin{figure}[H]
    \centering
    \includegraphics[width=0.8\textwidth]{screen-team-selection.png}
    \caption{Team Selection Screen}
    \label{fig:screen-teamselection}
\end{figure}

\subsection{Dashboard Screen}

\textit{Controller:} \texttt{DashboardController} \quad|\quad \textit{FXML:} \texttt{dashboard.fxml}

The Dashboard is the central hub of the game and the screen the user returns to after every
match (Requirement TM-2, LM-7, LM-9). It is divided into a top navigation bar, a left content
column, and a right standings sidebar.

\textbf{Top navigation bar} displays the managed team's name in amber and provides four navigation
buttons:

\begin{itemize}
    \item \textbf{Fixtures} $\rightarrow$ navigates to the Fixture screen
    \item \textbf{Squad} $\rightarrow$ navigates to the Squad screen
    \item \textbf{Table} $\rightarrow$ navigates to the Standings screen
    \item \textbf{Main Menu} $\rightarrow$ returns to main menu after a confirmation dialog
\end{itemize}

\textbf{Left column} contains three stacked cards:

\begin{itemize}
    \item \textit{Week/Season card} — displays the current matchday number (e.g.\ ``Matchday 5 / 38''), the team's current league position (e.g.\ ``3rd place''), and the season record (W/D/L, points, goal difference).

    \item \textit{Next Fixture card} — shows the upcoming match (e.g.\ ``Ironclad FC vs Silverbridge United''), the venue (Home/Away) and matchday number. The Play Match button triggers \texttt{SportManager.startMatch()}. A line beneath the button shows the lineup validity status — green if the starting XI is complete, amber warning if the lineup needs attention.

    \item \textit{Injury Report card} — lists every currently injured player in the squad with their position and the number of matches remaining in their recovery. If no players are injured, a green confirmation message is shown.
\end{itemize}

\textbf{Right sidebar} contains a mini standings table showing the top six clubs with position, name,
points, and goal difference. If the managed team sits below sixth place, their row is appended
below a continuation ellipsis. The managed team's row is always highlighted in amber.

When the season is complete (\texttt{isSeasonFinished()} returns true), the Dashboard replaces the
fixture card with a season-over banner displaying the champion's name.

\begin{figure}[H]
    \centering
    \includegraphics[width=0.9\textwidth]{screen-dashboard.png}
    \caption{Dashboard Screen}
    \label{fig:screen-dashboard}
\end{figure}

\subsection{Fixture Screen}

\textit{Controller:} \texttt{FixtureController} \quad|\quad \textit{FXML:} \texttt{fixture.fxml}

The Fixture screen displays the full season schedule across all matchdays (Requirement
TM-2). It is reached from the Dashboard's Fixtures button and calls
\texttt{SportManager.showFixtureData()} to retrieve all rounds.

The screen body contains a single TableView with five columns: \textbf{Wk} (matchday number),
\textbf{Home}, \textbf{Score}, \textbf{Away}, and \textbf{Status}. On load, the view automatically scrolls to the current
matchday so the user does not need to search manually.

Each row is styled based on its status:

\begin{itemize}
    \item \textit{Played} rows show the final score in the Score column and use muted grey text to indicate completed fixtures.
    \item \textit{Current matchday} rows are highlighted with an amber ``$\blacktriangleright$ NOW'' badge in the Status column.
    \item \textit{Upcoming} rows show ``vs'' in the Score column.
\end{itemize}

All rows where the managed team is involved (home or away) are highlighted with an amber
background tint, making it easy to track the team's schedule within the full fixture list. A Back
button in the top bar returns to the Dashboard.

\begin{figure}[H]
    \centering
    \includegraphics[width=0.9\textwidth]{screen-fixture.png}
    \caption{Fixture Screen}
    \label{fig:screen-fixture}
\end{figure}

\subsection{Squad Screen}

\textit{Controller:} \texttt{SquadController} \quad|\quad \textit{FXML:} \texttt{squad.fxml}

The Squad screen provides a complete view of the managed team's roster and coaching staff
(Requirements TM-1, LM-4, IM-2). It is reached from the Dashboard's Squad button.

The top bar displays the team name and an availability count (e.g.\ ``19 available / 22 total'').

The \textbf{Players table} occupies the majority of the screen and contains five columns:

\begin{itemize}
    \item \textbf{Name} — the player's full name
    \item \textbf{Position} — abbreviated position code (GK, CB, CM, ST, etc.)
    \item \textbf{Skill} — the player's overall skill rating (55–89)
    \item \textbf{Status} — ``FIT'' in green or ``INJ (N)'' in red, where N is the number of matches remaining
    \item \textbf{Attributes} — a compact summary of the player's sport-specific ratings (e.g.\ \texttt{Pace:78  Shooting:72  Passing:65})
\end{itemize}

Injured players have their entire row shaded in red, giving an immediate visual indication of
squad availability. The Status cell's text colour reinforces this: green for available, red for
injured.

The \textbf{Coaching Staff table} below the player table lists all coaches with their name, role (Head
Coach, Assistant, GK Coach, Fitness Coach), preferred formation, training skill rating, and
motivation skill rating.

A Back button returns to the Dashboard.

\begin{figure}[H]
    \centering
    \includegraphics[width=0.9\textwidth]{screen-squad.png}
    \caption{Squad Screen}
    \label{fig:screen-squad}
\end{figure}

\subsection{Match Screen}

\textit{Controller:} \texttt{MatchScreenController} \quad|\quad \textit{FXML:} \texttt{match-screen.fxml}

The Match Screen handles the full match experience: it shows the live score, renders match
events segment by segment, and provides the half-time intervention panel (Requirements
MS-1 through MS-5).

\textbf{Scoreboard panel} at the top displays the home team name on the left, the away team name
on the right, and the current score in large amber digits in the centre. A matchday and league
label appears below the score.

\textbf{Event log} occupies the left portion of the centre area inside a scrollable card. Events are
added incrementally as each segment is simulated. Each event type is colour-coded:

\begin{itemize}
    \item Goal events
    \item Injury events
    \item Yellow card events
    \item Substitution events
    \item Tactic change events
\end{itemize}

A segment header (e.g.\ ``FIRST HALF'') and a segment score summary (e.g.\ ``End of First
Half: 1\,--\,0'') are inserted between segments.

\textbf{Intervention panel} appears on the right side at half time and contains:

\begin{itemize}
    \item A substitution section with two dropdowns (remove player / add from bench) and a counter tracking substitutions used out of a maximum of three
    \item A tactic section with a dropdown showing all available formations, and an ``Apply Tactic'' button
\end{itemize}

The bottom bar holds the simulation button, whose label updates to reflect the next segment
(e.g.\ ``$\blacktriangleright$ Play First Half'' $\rightarrow$ ``$\blacktriangleright$ Play Second Half''). After the match ends, the simulation button is
replaced by a Continue button which calls \texttt{SportManager.advanceWeek()}, finalising the round
and returning to the Dashboard.

\begin{figure}[H]
    \centering
    \includegraphics[width=0.9\textwidth]{screen-match.png}
    \caption{Match Screen}
    \label{fig:screen-match}
\end{figure}

\subsection{League Table Screen}

\textit{Controller:} \texttt{LeagueTableController} \quad|\quad \textit{FXML:} \texttt{league-table.fxml}

The League Table screen shows the full sorted standings for the current season (Requirement
LM-7). It is accessible from the Dashboard's Table button and calls
\texttt{SportManager.showLeagueTableData()} which returns a \texttt{List<StandingEntry>} already sorted by
points $\rightarrow$ goal difference $\rightarrow$ goals for.

The screen body contains a TableView with ten columns: \textbf{\#} (position), \textbf{Club}, \textbf{P} (played), \textbf{W}, \textbf{D},
\textbf{L}, \textbf{GF}, \textbf{GA}, \textbf{GD}, and \textbf{Pts}. All numeric columns are centre-aligned; the Pts column is bold. The
managed team's row is highlighted in amber to allow the user to immediately identify their
standing among all twenty clubs.

The top bar shows the league name and the current matchday. A Back button returns to the
Dashboard.

\begin{figure}[H]
    \centering
    \includegraphics[width=0.9\textwidth]{screen-league-table.png}
    \caption{League Table Screen}
    \label{fig:screen-leaguetable}
\end{figure}

\subsection{Visual Design Principles}

All screens adhere to the following shared design principles:

\begin{itemize}
    \item \textbf{Colour palette:} Deep navy background (\texttt{\#0a0e1a}), card surfaces (\texttt{\#1c2235}), amber accent (\texttt{\#f59e0b}), green success (\texttt{\#10b981}), red danger (\texttt{\#ef4444}).

    \item \textbf{Typography:} Bold uppercase labels for section headers with increased letter-spacing; regular weight for data rows.

    \item \textbf{Consistent navigation:} Every secondary screen has a ``$\leftarrow$ Back'' button in the top-left corner returning to the Dashboard, creating a predictable navigation model.

    \item \textbf{State feedback:} Injured rows, managed-team rows, current-week rows, and played fixtures each have distinct visual treatments so the user can read game state at a glance without requiring tooltips or popups.

    \item \textbf{No sport-specific visuals in shared screens:} The League Table, Fixture screen, and Squad screen render correctly for any sport because they only read from the abstract \texttt{StandingEntry}, \texttt{Match}, and \texttt{Player} types — consistent with the UI independence requirement (NFR-02).
\end{itemize}

% ============================================================
%  11. EXTENSIBILITY
% ============================================================
\section{Extensibility}

\subsection{Design Principle}

The Sports Manager framework is built so that new sports can be added without modifying any
existing code in the UI or Core System layers. This follows the Open/Closed principle which
means the system is open for extension but closed for modification. Because no class in the
Core System or UI holds a reference to any concrete sport class, there is no existing code to
change when a new sport is introduced.

\subsection{How to Add a New Sport}

Adding a new sport requires only four steps:

\begin{enumerate}
    \item \textbf{Implement the Sport interface:} Create a class such as \texttt{VolleyballSport} that defines all sport-specific behavior: player generation, match simulation, points calculation, and injury logic.

    \item \textbf{Extend the abstract classes:} Create \texttt{VolleyballMatch}, \texttt{VolleyballPlayer}, \texttt{VolleyballTeam}, and \texttt{VolleyballLeague} by extending the corresponding abstract classes and adding sport-specific attributes and rules.

    \item \textbf{Register in SportFactory:} Add a single new case to \texttt{SportFactory.createSport()} so the factory returns the correct implementation when the user selects the sport.

    \item \textbf{Add to the selection screen:} Add the sport name to the list in \texttt{SportSelectionController}.
\end{enumerate}

\subsection{What Remains Unchanged}

When a new sport is added, every existing class remains untouched — all UI controllers,
SportManager, GameSession, League, Match, Team, Player, and all game logic. The only
existing file modified is SportFactory, where one line is added. This confirms that the
architecture fully satisfies the project requirement that the UI relies on an interface rather than
a real implementation.

% ============================================================
%  12. TECHNOLOGY STACK
% ============================================================
\section{Technology Stack}

The following technologies will be used:

\begin{itemize}
    \item \textbf{Language:} Java
    \item \textbf{GUI Framework:} JavaFX
    \item \textbf{Testing Framework:} JUnit 5
    \item \textbf{Build System:} Maven
    \item \textbf{Deployment:} \texttt{jpackage} to generate a Windows installer
    \item \textbf{Version Control:} Git
\end{itemize}

% ============================================================
\end{document}
